% GENERATED FILE. Edit canonical lessons, then run npm run generate:books.
% canonical-source-hash: fnv1a64:bc7fd2d5dbb6f81a
% canonical-lessons: PA-R158-food-quality-recall, PA-R158-truth-and-food-check

\chapter{Recall Familiar Quality Words at a Distance}
\label{ch:pa-quality-words-at-a-distance}
\hlchaptermodality{158}

\begin{chapteropening}
\textbf{By the end of this chapter:} \emph{I can retrieve five already-taught Punjabi quality words in familiar food and truth contexts without confusing this separate review with the scored A1 writing tasks.}

At a long interval, identify stale, ripe, unripe, true, and false from short familiar cues without introducing new Punjabi or borrowing the timed writing checkpoints.
\end{chapteropening}

\section[\texorpdfstring{\pa{ਬੇਹਾ} · \pa{ਪੱਕਾ} · \pa{ਕੱਚਾ} (behā · pakkā · kaccā)}{ਬੇਹਾ · ਪੱਕਾ · ਕੱਚਾ (behā · pakkā · kaccā)}]{Three old quality words at a food stall}

\label{lesson:PA-R158-food-quality-recall}



\begin{quote}
Say the old word for \textbf{stale}. (\textbf{\pa{ਬੇਹਾ}}, \emph{behā}.) Cover the answer before the next cue.
\end{quote}

\subsection*{Your turn}
Imagine familiar food at a stall. Say only the old Punjabi describing word; the English scenes do not add a Punjabi noun to learn.

\begin{itemize}
\raggedright
  \item Food left too long is \textbf{stale}. {[}PAUSE 4s{]} \textbf{\pa{ਬੇਹਾ}} (\emph{behā}).
  \item Fruit ready to eat is \textbf{ripe}. {[}PAUSE 4s{]} \textbf{\pa{ਪੱਕਾ}} (\emph{pakkā}).
  \item Fruit not yet ripe is \textbf{unripe}. {[}PAUSE 4s{]} \textbf{\pa{ਕੱਚਾ}} (\emph{kaccā}).
\end{itemize}

The words already carried these meanings in Chapter 146. \emph{Pakkā} can also mean firm; \emph{kaccā} can also mean raw. Do not add a new sense from this review.

\begin{tcolorbox}[breakable,colback=teal!4,colframe=teal!35!black,title={Before you move on}]
Cover the three answers. {[}PAUSE 4s each{]} Say the Punjabi for \textbf{unripe}, \textbf{stale}, then \textbf{ripe}. Check only after all three: \textbf{\pa{ਕੱਚਾ}}, \textbf{\pa{ਬੇਹਾ}}, \textbf{\pa{ਪੱਕਾ}}. If one slipped, revisit just that old Chapter 146 word.
\end{tcolorbox}
\section[\texorpdfstring{\pa{ਸੱਚਾ} · \pa{ਝੂਠਾ} (saccā · jhūṭhā)}{ਸੱਚਾ · ਝੂਠਾ (saccā · jhūṭhā)}]{Is it true? Return to all five}

\label{lesson:PA-R158-truth-and-food-check}



\begin{quote}
Say the word for \textbf{true}, then \textbf{false or lying}. Check: \textbf{\pa{ਸੱਚਾ}} (\emph{saccā}), then \textbf{\pa{ਝੂਠਾ}} (\emph{jhūṭhā}).
\end{quote}

\subsection*{Your turn}
An account that matches what happened is \textbf{true}: {[}PAUSE 3s{]} \textbf{\pa{ਸੱਚਾ}}. An account that does not is \textbf{false}: {[}PAUSE 3s{]} \textbf{\pa{ਝੂਠਾ}}. These English cues add no new Punjabi statement to memorise.

\begin{tcolorbox}[breakable,colback=teal!4,colframe=teal!35!black,title={Before you move on}]
Cover the earlier answers. {[}PAUSE 4s each{]} Say one old Punjabi word for each: \textbf{false}, \textbf{ripe}, \textbf{stale}, \textbf{true}, \textbf{unripe}.

Check in that same order: \textbf{\pa{ਝੂਠਾ}}, \textbf{\pa{ਪੱਕਾ}}, \textbf{\pa{ਬੇਹਾ}}, \textbf{\pa{ਸੱਚਾ}}, \textbf{\pa{ਕੱਚਾ}}. Reopen only a missed Chapter 146 word. These five distant retrievals do not turn the earlier short writing drills into an A1 mock.
\end{tcolorbox}
